% Problem record op_ebee7d5c442d81a4. This TeX file is derived from
% database/problems_json/op_ebee7d5c442d81a4.json by scripts/sync-tex.mjs; edit the JSON record
% and rerun the script rather than editing this file.
\section{Complete facet descriptions for Bell polytopes}
\label{sec:op_ebee7d5c442d81a4}

\paragraph{Problem.}
Determine complete facet descriptions for local-behavior polytopes beyond the
presently solved Bell scenarios, either for a specified unresolved finite
scenario or for a nontrivial infinite family with additional structure.  For
positive integers $N$, $M$, and $K$, let
$\mathbf{x}\in\{1,\ldots,M\}^{N}$ denote the measurement settings and
$\mathbf{a}\in\{1,\ldots,K\}^{N}$ the outcomes.  The relevant local polytope
is
\begin{equation}
  \mathcal{L}_{N,M,K}
  :=\operatorname{conv}\!\left\{
    p_{\mathbf f}:p_{\mathbf f}(\mathbf a\mid\mathbf x)
    =\prod_{i=1}^{N}\mathbf{1}\!\left\{a_i=f_i(x_i)\right\},\quad
    f_i:\{1,\ldots,M\}\to\{1,\ldots,K\}
  \right\}.
  \label{eq:p11-local-polytope}
\end{equation}
Here $\mathbf 1\{\cdot\}$ is the indicator function.
The task is to characterize all facet-defining inequalities of the polytope
in Eq.~\eqref{eq:p11-local-polytope} in a chosen unresolved regime, modulo
permutations of parties, settings, and outcomes and modulo liftings obtained
by adjoining redundant settings or outcomes.  A solution must carry a proof
of completeness, not merely generate a large collection of facets.

\subsection*{Status}

Unsolved

\subsection*{Source}

Kr\"uger and Werner explicitly pose the search for completeness-certified
facet descriptions of Bell polytopes in tractable finite scenarios and
structured families \sourcecite{ref:p11-krueger-werner}{KW05}.

\subsection*{Progress}

\begin{itemize}
  \item Fine proved that positivity, normalization, and the
  Clauser--Horne--Shimony--Holt inequalities completely describe the
  $(N,M,K)=(2,2,2)$ local scenario.  This settles only the smallest
  nontrivial case
  \sourcecite{ref:p11-fine}{Fin82}.

  \item Werner and Wolf classified all full-correlation facets for arbitrary
  $N$ when every party has two dichotomic observables.  Their theorem concerns
  a correlator projection of Eq.~\eqref{eq:p11-local-polytope}, not the full
  conditional-probability polytope
  \sourcecite{ref:p11-werner-wolf}{WW01}.

  \item Staufenbiel's polyhedral-sampling method found more than
  $1.29\times10^8$ inequivalent facet classes in one unresolved scenario, but
  the method intentionally sacrifices completeness.  Thus the computation
  enlarges the known catalog without supplying the required certificate
  \sourcecite{ref:p11-staufenbiel}{Sta26}.
\end{itemize}

\subsection*{References}

\begin{enumerate}[label={},leftmargin=4.8em,labelsep=0.5em]
  \footnotesize
  \item[\textup{[Fin82]}]\label{ref:p11-fine}
  A. Fine,
  ``Hidden Variables, Joint Probability, and the Bell Inequalities,''
  \emph{Physical Review Letters} \textbf{48}, 291--295 (1982).
  \href{https://doi.org/10.1103/PhysRevLett.48.291}{doi:10.1103/PhysRevLett.48.291}.

  \item[\textup{[WW01]}]\label{ref:p11-werner-wolf}
  R. F. Werner and M. M. Wolf,
  ``All Multipartite Bell Correlation Inequalities for Two Dichotomic
  Observables per Site,''
  \emph{Physical Review A} \textbf{64}, 032112 (2001).
  \href{https://doi.org/10.1103/PhysRevA.64.032112}{doi:10.1103/PhysRevA.64.032112};
  \href{https://arxiv.org/abs/quant-ph/0102024}{arXiv:quant-ph/0102024}.

  \item[\textup{[Sta26]}]\label{ref:p11-staufenbiel}
  C. Staufenbiel,
  ``Bell Inequalities from Polyhedral Sampling,'' arXiv:2604.22859 (2026).\newline
  \href{https://doi.org/10.48550/arXiv.2604.22859}{doi:10.48550/arXiv.2604.22859};
  \href{https://arxiv.org/abs/2604.22859}{arXiv:2604.22859}.

  \item[\textup{[KW05]}]\label{ref:p11-krueger-werner}
  O. Krüger and R. F. Werner (eds.),
  ``Some Open Problems in Quantum Information Theory,''
  arXiv:quant-ph/0504166 (2005).
  \href{https://doi.org/10.48550/arXiv.quant-ph/0504166}{doi:10.48550/arXiv.quant-ph/0504166};
  \href{https://arxiv.org/abs/quant-ph/0504166}{arXiv:quant-ph/0504166}.
\end{enumerate}

\subsection*{Comment}

The source problem explicitly replaces an implausible classification uniform
in all $(N,M,K)$ by completeness-certified finite cases and structured
families \sourcecite{ref:p11-krueger-werner}{KW05}.  The unresolved gap is a
new theorem of completeness for one such nontrivial regime; heuristic
completeness or sampling without a certificate is insufficient.

\subsection*{Field}

Quantum Resource Theory

\subsection*{Topic}

Bell nonlocality

\subsection*{ID}

\texttt{op\_ebee7d5c442d81a4}
